\subsection{An isolated dyadic wave and an actual-coefficient index test}

The simple poles of P-F31 permit a strict-index test using a specified
cyclotomic trace and an elementary bound on an actual coefficient.
The reduction is reproduced with its complete proof.  The conditions
explicitly stated as implications here are discharged by PNEW-005 below.

\begin{proposition}[P-F56: maximal two-part poles and paired constituents]
\label{prop:cubicpairedwave}
Use the ordinary cubic multiplicity algebra $T^{(n)}$ and its
Hilbert function $f_n$ from P-F31, with $n\ge4$, and put
\[
 M=2^{n-1},\qquad g=\gcd(n,3),\qquad d_0=n/g,
 \qquad m=\binom{n+2}{3}.
\]
Every pole of $H_n(t)$ is a root of unity whose order has
two-part at most $M$. Among roots whose order has two-part exactly
$M$, the only poles are the primitive $M$th roots, and each of
those poles is simple. More precisely, if $\eta$ is a root of
unity of order $q$ with $v_2(q)\ge n-1$, then
\[
 \eta\text{ is a pole of }H_n
       \quad\Longleftrightarrow\quad q=M;
 \qquad q=M\Longrightarrow\text{the pole is simple}.
 \tag{P-F56a}
\]

Let $L$ be the least quasipolynomial period of $f_n$, and write
$f_n(s)=q_r(s)$ for $s\equiv r\pmod L$, with
$0\le r<L$. Then
\[
                     v_2(L)=n-1.                            \tag{P-F56b}
\]
Fix $\zeta=e^{2\pi i/M}$ and use the following actual-residue normalization:
\[
 \mathcal C_n(\zeta)=\lim_{t\to\zeta}(1-t^M)H_n(t),
 \qquad A_n(\zeta)=\mathcal C_n(\zeta)/M.
\]
For each integer $r$ define the rational, $M$-periodic wave
\[
 \begin{split}
 W_n(r)
   &=\sum_{\substack{1\le a<M\\a\text{ odd}}}
                      A_n(\zeta^a)\zeta^{-ar}\\
   &=\frac1M\operatorname{Tr}_{\mathbb Q(\zeta)/\mathbb Q}
                  \bigl(\mathcal C_n(\zeta)\zeta^{-r}\bigr).
 \end{split}                                                 \tag{P-F56c}
\]
For every residue $r$ and every complex argument $X$ one has
the exact polynomial identity
\[
               q_r(X)-q_{r+L/2}(X)=2W_n(r),                 \tag{P-F56d}
\]
where subscripts are read modulo $L$. In particular,
\[
 \begin{split}
 f_n\bigl(p(f_n)\bigr)&\ge2W_n\bigl(p(f_n)\bigr),\\
 f_n\bigl(s(f_n)+1\bigr)&\ge2W_n\bigl(s(f_n)+1\bigr).
 \end{split}                                                 \tag{P-F56e}
\]
These are necessary conditions for the literal indices of
GCT VI, Definitions~1.2.2 and 1.2.4, printed pp.~13--14.
They require no bound on the average polynomial $B_n$ or on
Fourier terms of smaller two-part.
More generally, the inequality $f_n(x)\ge2W_n(x)$ holds
at every integer $x\ge p(f_n)$ and at every integer
$x\ge s(f_n)+1$. Thus a single integer $x\ge0$ satisfying
\[
          2W_n(x)>f_n(x)
       \quad\Longrightarrow\quad p(f_n)>x,
                                  \quad s(f_n)\ge x.
 \tag{P-F56e'}
\]

For every fixed real $K\ge0$ and every integer $0\le x\le n^K$,
\[
 f_n(x)\le\binom{m+d_0x-1}{m-1}
     \le\bigl(2n^{\max\{3,K+1\}}\bigr)^{n^3}.
 \tag{P-F56f}
\]
Consequently the following would suffice to refute both narrow
polynomial-index hypotheses on this cubic family: there exist
$\epsilon>0$ and an unbounded set $\mathcal N$ of dimensions
such that, for every fixed $K\ge0$ and every sufficiently large
$n\in\mathcal N$, there is an integer $x$ satisfying
\[
 \lceil n^K\rceil\le x\le\lceil n^K\rceil+2,
                         \qquad W_n(x)\ge2^{\epsilon n^4}.
 \tag{P-F56g}
\]
Under this additional condition, both $p(f_n)$ and $s(f_n)$
exceed every fixed polynomial in $n$ along $\mathcal N$.
PNEW-005 below supplies a witness in a window of width at most nine,
which is sufficient for the same conclusion. In the present conditional criterion the wave can have
either sign a priori, and the absolute value of one Fourier
coefficient alone does not establish (P-F56g).
The window of three consecutive integers is a sufficient
target, not a proved oscillation pattern; any bounded-width
window, or suitable polynomially bounded witnesses tending
beyond every fixed polynomial threshold, would work as well.
\end{proposition}

\begin{proof}
We first prove a fixed-point criterion for absence of a Hilbert
pole, including possible nonreduced structure at the origin.
Let $R$ be any finitely generated positively graded complex
algebra with $R_0=\mathbb C$, and let $\eta$ have order $q$.
The ideal defining its grading fixed scheme is
\[
                    I_q=(R_d:q\nmid d).                     \tag{P-F56h}
\]
Indeed a homogeneous element of degree $d$ changes by
the scalar $\eta^d$, so its difference from its transform
is zero exactly when $q\mid d$.
Choose homogeneous ideal generators $a_1,\ldots,a_l$ of
$I_q$, each of degree $e_i$ not divisible by $q$.
Such a finite choice is possible by Noetherianity, starting
with the indicated homogeneous generating set.
If the fixed scheme is supported only at the graded origin,
then $\sqrt{I_q}=R_{>0}$ and $R/I_q$ is a finite-dimensional
complex algebra. To see this directly, a power of the finitely
generated ideal $R_{>0}$ lies in $I_q$, and only finitely many
monomials of bounded length in homogeneous algebra generators
survive.

The homology modules $N_j$ of the graded Koszul complex
$K_R(a_1,\ldots,a_l)$ are finite $R$-modules annihilated by
$I_q$: multiplication by each $a_i$ is null-homotopic on
the Koszul complex. Thus all $N_j$ are finite-dimensional.
The degreewise Euler identity gives
\[
 H_R(t)\prod_{i=1}^l(1-t^{e_i})
                    =\sum_j(-1)^j\operatorname{Hilb}_{N_j}(t).
 \tag{P-F56i}
\]
The right side is a polynomial and every factor on the left
is nonzero at $t=\eta$. Hence $H_R(t)$ is regular at $\eta$.
This is the zero-dimensional case of the Koszul calculation
in P-F23; reducedness of the fixed scheme is unnecessary.

Apply this criterion to $R=T^{(n)}$ and $X=\operatorname{Spec}R$.
Let $q=2^ab$ with $b$ odd and $a\ge n-1$.
Suppose a nonzero closed point $p\in X$ is fixed by a primitive
$q$th root $\eta$. It is then fixed by
$\eta_2=\eta^b$, which has order $2^a$.
We need the closed-orbit lift for this pure two-primary
subgroup, with no assumption about the odd part of $q$.
Here are the details extending the first part of P-F26.

Put $\gamma=3/g$, an odd integer. The canonical grading
character on quotient points is
$(\det B)^{-\gamma}$ for $B\in GL_n$, as in (P-F26b).
There is a primitive $2^a$th root $\theta$ such that
$\theta^{-\gamma}=\eta_2$.
Choose $B\in GL_n$ with $\det B=\theta$, and choose
a cubic $F$ in the unique closed $SL_n$ orbit over $p$.
Both $F$ and $B\cdot F$ have closed orbits and the same
quotient point. Therefore $B\cdot F=S\cdot F$ for some
$S\in SL_n$, and $C=S^{-1}B$ fixes $F$ exactly with
$\det C=\theta$.
The unique closed orbit in each affine reductive quotient
fiber is the standard input already used in P-F26; an exact
primary locator is Luna, \emph{Slices \'etales},
Section~I.1, printed p.~84:
\url{https://www.numdam.org/item/10.24033/msmf.110.pdf}.

The semisimple part $C_s$ also fixes $F$ and has determinant
$\theta$. Indeed in the commuting Jordan decomposition,
an eigenspace of $C_s$ with eigenvalue different from one
cannot contain a fixed vector of $C_sC_u$, because the
unipotent part $C_u$ has all eigenvalues one.
Diagonalize $C_s$ and let $\Gamma\subset\mathbb C^*$ be
the finitely generated group generated by its eigenvalues.
Project $\Gamma$ to its finite two-primary torsion subgroup,
using a decomposition into its finite cyclic torsion subgroup
and a free complement. This projection fixes $\theta$.
Applying it to every eigenvalue preserves every monomial
relation expressing $C_sF=F$.
The resulting diagonal matrix $A$ fixes $F$, has finite
two-power order, and has determinant of order exactly $2^a$.

Since $p$ is nonzero, a positive-degree invariant is nonzero
at $F$. In any eigenbasis, one of its nonzero coefficient
monomials uses only monomials supported in $F$.
Invariance under the diagonal $SL_n$ torus forces the
corresponding nonzero nonnegative integral combination of supported
cubic exponent vectors to be a positive rational multiple
of $\mathbf1$. Thus $\mathbf1$ belongs to the rational
span of the support, precisely the hypothesis of P-F30.
Theorem P-F30 gives
\[
                       2^a=\operatorname{ord}(\det A)\le M.
 \tag{P-F56j}
\]
If $a>n-1$ this is a contradiction, so no nonzero fixed
point exists. If $a=n-1$, equality in P-F30 makes $F$
equivalent to the chain cubic, hence $p$ lies on its
nonzero grading orbit. The exact stabilizer of every
point of that orbit is $\mu_M$ by P-F28, including the
converse lift argument in (P-F28d).
The original root $\eta$ can fix $p$ only if $\eta\in\mu_M$.
Since its order is $q=Mb$, this requires $b=1$.

It follows that for every $q$ in (P-F56a) except $q=M$
the fixed scheme has no nonzero closed point. Its support
is consequently just the graded origin, because a closed
subset of an affine scheme of finite type over $\mathbb C$
is determined by its closed points. Equations
(P-F56h)--(P-F56i) show that there is no Hilbert pole there.
For $q=M$, P-F31 proves an actual simple pole at every
primitive root. This proves all of (P-F56a).

For completeness, the least quasipolynomial period is the
least common multiple of the orders of the Hilbert poles.
Indeed the finite Fourier transforms are
\[
 P_\rho(X)=\frac1L\sum_{r=0}^{L-1}\rho^r q_r(X),\qquad
 q_r(X)=\sum_{\rho^L=1}\rho^{-r}P_\rho(X).                 \tag{P-F56k}
\]
The generating function of a nonzero polynomial times
$\rho^{-s}$ has its only pole at $t=\rho$, with order
one more than its polynomial degree. Thus $P_\rho\ne0$
if and only if $\rho$ is a pole of $H_n$.
Let $L_0$ be the least common multiple of those root orders.
The second expression in (P-F56k) is periodic in $r$
with period $L_0$, so it exhibits $L_0$ as a quasipolynomial
period. Conversely $L_0\mid L$. Minimality gives $L=L_0$.
Now (P-F56a) and the primitive $M$th poles give (P-F56b).

In particular $L/M$ is odd. If $\rho$ is any Hilbert pole
other than a primitive $M$th root, its order has two-part
at most $M/2$ and divides $L/2$. Its contribution in the
difference of (P-F56k) at $r$ and $r+L/2$ therefore
vanishes. For $\rho=\zeta^a$ with $a$ odd,
$\rho^{-L/2}=-1$.
At a simple pole $\rho$, the term with its principal part is
$A_n(\rho)/(1-t/\rho)$: indeed
$\lim_{t\to\rho}(1-t^M)/(1-t/\rho)=M$.
Its coefficient at degree $s$ is $A_n(\rho)\rho^{-s}$,
so $P_{\zeta^a}(X)=A_n(\zeta^a)$, independent of $X$.
Subtracting the two Fourier formulas proves (P-F56d).
It also shows that every nonconstant coefficient of the
constituent polynomials has period dividing $L/2$.
The rational Hilbert series is defined over $\mathbb Q$,
so Galois conjugation sends $\mathcal C_n(\zeta)$ to
$\mathcal C_n(\zeta^a)$. This proves the trace expression
and rationality in (P-F56c). Its periodicity is immediate;
also $W_n(r+M/2)=-W_n(r)$.

After the coefficient-positivity shift $p=p(f_n)$, all
constituents evaluated at $p$ are nonnegative, since these
values are their shifted constant coefficients.
Take the residue $r\equiv p\pmod L$ in (P-F56d).
Its first term at $X=p$ is the actual value $f_n(p)$,
and its second term is nonnegative.
After a strict-saturation shift $s=s(f_n)$, every nonzero
constituent at $s+1$ is positive, and every identically
zero constituent has value zero. Taking
$r\equiv s+1\pmod L$ proves the second inequality of
(P-F56e). This uses exactly the source conventions, not
positivity only on each constituent's own progression.
These tests persist at every later integer argument.
For positivity, each polynomial $q_r(p+X)$ has nonnegative
coefficients and hence is nonnegative for all real $X\ge0$.
For saturation, the definition directly gives
$q_r(s+j)\ge0$ for every integer $j\ge1$.
Applying (P-F56d) at the residue of any such later argument
proves (P-F56e'). No derivative positivity is being inferred
from strict saturation.

Finally the multiplicity $f_n(x)$ is bounded by the
dimension of the ambient symmetric power of an
$m$-dimensional space. Since $m\le n^3$ and
$d_0x\le n^{K+1}$,
\[
 \binom{m+d_0x-1}{m-1}
      \le(m+d_0x)^{m-1}
      \le\bigl(2n^{\max\{3,K+1\}}\bigr)^{n^3},
\]
including $x=0$. This proves (P-F56f).
Its base-two logarithm is $O_K(n^3\log n)$.
The witness in (P-F56g) is at most $n^{K+1}$ for all
sufficiently large $n$. Apply (P-F56f) with $K+1$.
Its actual coefficient is smaller than $2W_n(x)$ for
sufficiently large $n\in\mathcal N$, so (P-F56e') gives
$p(f_n)>x\ge n^K$ and $s(f_n)\ge x\ge n^K$.
Applying this with a larger fixed exponent proves that
neither index is bounded by any polynomial in $n$.
The GCT parameter is $h=m^2+1+n=\Theta(n^6)$ by P-F31,
so a polynomial bound in that parameter would also fail.
This establishes the conditional criterion. The analytic estimates
and bounded-window witness are proved in PNEW-005 below.\hfill$\square$
\end{proof}


\paragraph{An exact positive multiset for the chain normal representation.}
There is also a numerator-free product specifying every residue
in the remaining estimate. Regard
$w_0,\ldots,w_{n-1}$ as residues modulo $M$, with the
$w_i=(-2)^i$ for $i<n-1$ and $w_{n-1}=0$, and form the multiset
\[
 \begin{gathered}
 \mathcal E_n=
 \{\{0,1,k\}:2\le k<n\}\cup\{\{0,2,3\}\},\\
 \mathcal S_n=\{w_i:0\le i\le n-2\}\ \sqcup\
 \{w_i+w_j+w_k:0\le i<j<k<n,\
       \{i,j,k\}\notin\mathcal E_n\}.
 \end{gathered}
 \tag{P-F56u}
\]
Repeated weights are retained. Every weight is nonzero modulo
$M$, and $|\mathcal S_n|=D=\binom n3$.
With the actual Hilbert-residue convention, the exact identity is
\[
 \mathcal C_n(\zeta)=
                 \prod_{b\in\mathcal S_n}(1-\zeta^{gb})^{-1}.
 \tag{P-F56v}
\]
This is an equality of complex numbers, including their phases,
not only an absolute-value identity.

To prove it, put $\eta=\zeta^g$ and let
$\alpha_i=\eta^{w_i}$ be the natural eigenvalues of the diagonal
chain automorphism. Write $V$ for their representation.
Its action on polynomial coordinate functions is the usual
$GL_n$ action on $\operatorname{Sym}^3 V$.
On an invariant of coefficient degree $d_0s$ the character is
$(\det)^{\gamma s}$, where $\gamma=3/g$.
Since $3\sum_i w_i\equiv1\pmod M$, one has
$(\prod_i\alpha_i)^\gamma=\zeta$.
Thus its action on the canonical multiplicity algebra is precisely
the ring automorphism multiplying degree $s$ by $\zeta^s$.
This choice fixes the convention required in (P-F56v).

The smooth quotient tangent sequence in P-F24, dualized to
coordinate functions, shows that the moving conormal character is
\[
 \operatorname{Sym}^3 V-V\otimes V^*.
 \tag{P-F56w}
\]
Indeed the removed orbit cotangent character is the adjoint
character $V\otimes V^*-\mathbf1$, and the remaining fixed
grading direction contributes the other $\mathbf1$.
This is an identity in the representation ring of the finite
cyclic group, with its left side realized by the actual
conormal representation.

Let $h_j$ and $e_j$ be the complete and elementary symmetric
characters in the $\alpha_i$.
The identity
$\prod_i(1+t\alpha_i)=
 \prod_i(1-t^2\alpha_i^2)/\prod_i(1-t\alpha_i)$
gives $h_3=e_3+h_1\sum_i\alpha_i^2$.
The multisets $\{\alpha_i^2\}$ and $\{\alpha_i^{-1}\}$
differ by replacing $\alpha_0^{-1}$ with one: this follows
from $\alpha_i^2=\alpha_{i+1}^{-1}$ for $i<n-1$
and $\alpha_{n-1}=1$.
Consequently (P-F56w) is
\[
                       e_3+h_1-\alpha_0^{-1}h_1.
 \tag{P-F56x}
\]
The trivial weight in $h_1$ cancels the term with $i=0$
in $\alpha_0^{-1}h_1$.
For every $2\le k<n$, the remaining subtracted weight
$w_k-w_0=w_k-1$ equals the weight of the triple
$\{0,1,k\}$, because $w_0+w_1=-1$.
The last subtracted weight is $w_1-w_0=-3$.
It equals the weight of $\{0,2,3\}$:
for $n\ge5$ the integer sum is $1+4-8=-3$;
for $n=4$ it is $1+4+0=5\equiv-3\pmod8$.
These removed triples are distinct.
The remaining character is exactly the positive multiset
(P-F56u). The count is
$(n-1)+\binom n3-(n-2)-1=D$.
Nonvanishing of its singleton weights is immediate, and
a distinct triple sum is nonzero by its unique smallest
two-adic valuation among the nonzero summands.
The residue formula (P-F31b) now gives (P-F56v).

For reference, put
$a_0=g(n+1)(n+2)/6\in\mathbb Z$.
The determinant of $\operatorname{Sym}^3 V$ is
$(\det V)^{\binom{n+2}2}$, whereas
$\det(V\otimes V^*)=1$.
It follows from (P-F56w) and $3\sum_i w_i\equiv1\pmod M$
that the sum of the exponents $gb$ in (P-F56v) is congruent
to $a_0$ modulo $M$.
Inverting each factor therefore gives the exact relations
\[
 \mathcal C_n(\zeta^{-1})
       =(-1)^D\zeta^{a_0}\mathcal C_n(\zeta),
 \qquad
 W_n(s)=(-1)^D W_n(-s-a_0).
 \tag{P-F56y}
\]
They fix a reciprocity symmetry, but do not prescribe the sign
of the wave at a specified nonnegative integer.

Finally, if $h=M/2$ and
$\mathcal C_n(\zeta)=\sum_{j=0}^{h-1}c_{n,j}\zeta^j$
is its unique representative in
$\mathbb Q[t]/(t^h+1)$, then for every integer $s$,
\[
             2W_n(s)=
       (-1)^{\lfloor s/h\rfloor}c_{n,\,s\bmod h}.
 \tag{P-F56z}
\]
Indeed the primitive dyadic trace of $\zeta^b$ is zero unless
$h\mid b$, and is $h(-1)^{b/h}$ otherwise. This follows
directly by summing the geometric series over the odd
exponents modulo $M$.
Only $j=s\bmod h$ survives in (P-F56c), giving (P-F56z).
In particular, at $0\le x<h$ the exact one-witness
obstruction is $f_n(x)<c_{n,x}$.
This is a coefficient reformulation of the unresolved signed
wave question, with no positivity of those coefficients assumed.

\paragraph{A contraction gate involving only quadratic-many factors.}
The unresolved mass comparison can be reduced further to a
dimension-increment inequality. This reduction proves the error
from the two selected sibling branches; it does not assume or
prove the remaining uniform contraction inequality.

For $n\ge4$ write $M_n=2^{n-1}$, let
$b_n=3^{-1}\pmod {M_n}$, and use the positive multiset
$\mathcal S_n$ of (P-F56u). For every odd residue $a\bmod M_n$
define
\[
 P_n(a)=
 \prod_{b\in\mathcal S_n}
       |1-\exp(2\pi iab/M_n)|^{-1},\qquad
 Q_n=\sum_{\substack{a\bmod M_n,\ a\ {\rm odd}\\
                               a\ne b_n,-b_n}}
                              \frac{P_n(a)}{P_n(b_n)}.
 \tag{P-F56aa}
\]
Every factor is nonzero. Multiplication of the primitive-root
index by the odd number $g=\gcd(n,3)$ permutes the odd residues.
Thus (P-F56v) implies exactly $Q_n=E_n/|Z_n|$; the common
normalization $1/M_n$ cancels.

For $n\ge5$, form the multiset of residues modulo $M_n$
\[
 \begin{split}
 \mathcal T_n={}&\{1\}\\
 &\sqcup\{1+w_j+w_k:2\le j<k<n,\ (j,k)\ne(2,3)\}\\
 &\sqcup\{2+w_k:3\le k<n\}\sqcup\{6\},
 \qquad |\mathcal T_n|=\binom{n-1}2,
 \end{split}
 \tag{P-F56ab}
\]
where $w_i=(-2)^i$ for $i<n-1$ and $w_{n-1}=0$,
as before. Define the positive increment product
$J_n(a)=\prod_{b\in\mathcal T_n}
 |1-\exp(2\pi iab/M_n)|^{-1}$.
The exact recursion is
\[
             P_n(a)=J_n(a)P_{n-1}(-a\bmod M_{n-1}).
 \tag{P-F56ac}
\]

Here is a sufficient uniform contraction hypothesis, involving
only the $\binom{n-1}2=O(n^2)$ factors of each increment:
there exist a fixed $0<\rho<1$ and an integer $n_0\ge5$ such that,
for every $n\ge n_0$ and every odd $a\bmod M_n$ with
$a\bmod M_{n-1}\notin\{b_{n-1},-b_{n-1}\}$,
\[
 J_n(a)+J_n(a+M_{n-1})\le \rho J_n(b_n).
 \tag{P-F56ad}
\]
PNEW-002--005 below prove this inequality with explicit constants.
We first show that it implies $Q_n\longrightarrow0$. The phase-window
argument at the end of this section then produces the positive
wave values used in PNEW-005.
In this implication there is no unproved finite base
certificate: any finite starting value of $Q_n$ suffices.

The selected sibling error is bounded unconditionally:
\[
 \delta_n:=
 \frac{P_n(b_n+M_{n-1})}{P_n(b_n)}
 \le
 2^{-n^3/6+4n^2-107n/6+19},
 \qquad
 \delta_n\le2^{-n^3/12}\quad(n\ge48).
 \tag{P-F56ae}
\]
The exact implication of (P-F56ad) is
\[
                         Q_n\le\rho Q_{n-1}+2\delta_n.
 \tag{P-F56af}
\]

\emph{Proof.}
Remove the first variable from the chain weights.
The remaining weights $w_1,\ldots,w_{n-1}$ are $-2$
times the weights of the $(n-1)$-variable chain, in the
same order. In the multiset (P-F56u), all triples with
smallest index at least one occur. The shifted smaller
multiset removes from these exactly the triples
$\{1,2,k\}$ for $3\le k<n$ and $\{1,3,4\}$.
The additional triples with smallest index zero are
$\{0,j,k\}$ for $2\le j<k<n$, except $\{0,2,3\}$.
The additional singleton is $w_0=1$.
These extra weights are (P-F56ab): the first removed
family has weight $w_1+w_2+w_k=2+w_k$, and the last
triple has weight six modulo $M_n$.
For $n\ge6$ its integer weight is $-2-8+16=6$;
for $n=5$ it is $-2-8+0=-10\equiv6\pmod {16}$.
Thus the multiset decomposition is exact, including
multiplicities and the smallest allowed dimension.
Its cardinality is
$1+(\binom{n-2}2-1)+(n-3)+1=\binom{n-1}2$.
All its factors are nonzero because it is a submultiset
of $\mathcal S_n$. Substitution of the root
$\exp(-4\pi ia/M_n)$ in the smaller product proves
(P-F56ac).

Put $M=M_n$ and $\eta=\exp(2\pi ib_n/M)$.
The two roots indexed by $b_n$ and $b_n+M/2$
have the same smaller-dimensional factor in (P-F56ac).
In the increment, every even weight gives the same
factor at $\eta$ and $-\eta$.
The odd weights of $\mathcal T_n$ are exactly
\[
 \mathcal O_n=
 \{1\}\sqcup
 \{1+w_j+w_k:2\le j<k<n,\ (j,k)\ne(2,3)\}.
 \]
It follows that
\[
                 \delta_n
       =\prod_{b\in\mathcal O_n}
                     \left|\tan\frac{\pi b_n b}{M}\right|.
 \tag{P-F56ag}
\]
No tangent is singular, since $b_n$ and every $b$ here
are odd and $M\ge16$.

Among these odd weights, the bulk pairs
$2\le j<k\le n-2$, excluding $(2,3)$,
give
$B=\binom{n-3}2-1$ weights divisible by three.
For such a weight $b=1+w_j+w_k$, one has
$|b|<2^{k+1}$ and $|b|\le M$.
Since $3b_n\equiv1\pmod M$,
$\tan(\pi b_n b/M)=\tan(\pi b/(3M))$.
On $[-\pi/3,\pi/3]$, the bounds
$|\sin u|\le|u|$ and $\cos u\ge1/2$ give
\[
 \left|\tan\frac{\pi b_n b}{M}\right|
       \le\frac{2\pi |b|}{3M}
       \le\frac{8\,2^k}{M}.
 \tag{P-F56ah}
\]
The sum of $n-1-k$ over these bulk pairs is
$\binom{n-2}3-(n-4)$: before the one excluded pair,
it counts triples $2\le j<k<\ell\le n-1$.
There are $n-2$ exceptional odd factors, namely the
singleton one and the pairs using the zero-weight
variable. Each is at most $M/2$.
Indeed an odd residue modulo $M$ is at distance at least
one from $M/2$, so its cosine in (P-F56ag) has absolute
value at least $\sin(\pi/M)\ge2/M$, while its sine has
absolute value at most one.
Combining these estimates gives
\[
 \begin{split}
 \log_2\delta_n
 &\le 3\left(\binom{n-3}2-1\right)
        -\binom{n-2}3+(n-4)+(n-2)^2\\
 &=-n^3/6+4n^2-107n/6+19 .
 \end{split}
 \tag{P-F56ai}
\]
For $n\ge48$, $4n^2\le n^3/12$ and
$-107n/6+19<0$, proving the last assertion of
(P-F56ae).

Finally group the odd residues modulo $M_n$ into their
two lifts above an odd residue modulo $M_{n-1}$.
Reduction of $b_n$ is $b_{n-1}$.
The two selected parent classes therefore have the two
selected children $\pm b_n$ and the two unselected
children $\pm b_n+M_{n-1}$.
Complex conjugation makes the latter two contributions
equal, with total normalized mass $2\delta_n$.
For every other parent, (P-F56ac) and (P-F56ad)
bound its two children by $\rho$ times its normalized
mass in $Q_{n-1}$.
Here
$P_n(b_n)=J_n(b_n)P_{n-1}(b_{n-1})$,
using conjugation for the minus sign in (P-F56ac).
This proves (P-F56af), with no factor of two lost
in the nonselected branches.

Since $\delta_n\to0$ and $\rho<1$, the recursion forces
$Q_n\to0$. Explicitly, for any $\epsilon>0$, choose
$N$ so that $2\delta_n\le\epsilon(1-\rho)/2$
for all $n>N$. Iterating gives
\[
 Q_n\le \rho^{n-N}Q_N+\epsilon/2
                  \quad(n\ge N),
 \tag{P-F56aj}
\]
and the first term tends to zero.
The starting value $Q_N$ is finite by its definition.
This completes the conditional contraction reduction
and the unconditional selected-sibling bound.



\paragraph{The bounded phase window, with its normalization.}
Put $r=3g\in\{3,9\}$, choose $a_*r\equiv1\pmod M$, and set
$Z_n=A_n(\zeta^{a_*})$ and
$E_n=\sum_{a\ {
m odd},\,a\ne\pm a_*}|A_n(\zeta^a)|$.
Then $Q_n=E_n/|Z_n|$ by (P-F56v) and (P-F56aa).
Write $ra_*=1+kM$.  Here $1\le k<r$ and $\gcd(k,r)=1$.
For any integer $X$ the $r$ numbers
$Z_n\zeta^{-a_*X}\exp(-2\pi i k j/r)$, $0\le j<r$,
are equally spaced on the circle of radius $|Z_n|$.
One has real part at least $\cos(\pi/r)|Z_n|\ge |Z_n|/2$.
Replacing that phase by $\zeta^{-a_*j}$ changes its real part
by at most $2\pi|Z_n|/M\le |Z_n|/4$ when $n\ge6$.
The selected conjugate pair therefore contributes at least
$|Z_n|/2$ at some $x\in[X,X+r-1]$.
If $Q_n\le1/4$, the remaining terms have absolute sum at most
$|Z_n|/4$, and hence $W_n(x)\ge |Z_n|/4$.
PNEW-004 below gives the required selected-amplitude lower bound
directly from the exact positive residue product.
